%% ============================================================
%%  Handout — Aula 2: Tomando boas decisões dado um modelo do mundo
%%  Curso: Aprendizagem por Reforço — UFU
%%  Prof. Marcelo Keese Albertini
%%  Adaptado e expandido de CS234 (Stanford), Emma Brunskill
%%  COMPILAR COM XeLaTeX (duas vezes)
%% ============================================================
\documentclass[11pt, a4paper]{article}
%% .sty do projeto na raiz do repositório (../)
\makeatletter
\def\input@path{{./}{../}}
\makeatother


\usepackage{handoutAprendizadoRL}
\usepackage{babel}
\babelprovide[import, main]{brazilian}

\setaula{Aula 2 — Decisões com modelo}
\setprof{Prof. Marcelo Keese Albertini}

%% Macros locais (mesmas da apresentação)
\newcommand{\Bop}{\mathcal{B}}
\newcommand{\Bpol}{\mathcal{B}^{\pol}}
\newcommand{\Rmax}{R_{\max}}

\begin{document}

\capa{Aula 2: Tomando sequências de boas decisões dado um modelo do mundo}
     {Processos de decisão de Markov \textbullet\ Avaliação de política \textbullet\ Iteração de política e de valor}
     {Prof. Marcelo Keese Albertini}
     {Adaptado de Emma Brunskill, CS234: Reinforcement Learning, Stanford University \textbullet\ 2026}

\paragraph{Contexto.} Esta é a única aula do curso em que o agente conhece \emph{tudo}
sobre o mundo: a dinâmica $P$ e a recompensa $\rec$. O problema deixa de ser de
aprendizagem e passa a ser de \textbf{planejamento} --- e é exatamente por isso que ela
fornece o alicerce teórico de todo o resto. Dado o modelo $(\gS, \gA, P, \rec, \gamma)$,
dois problemas nos ocupam:

\begin{enumerate}
  \item \textbf{Predição (avaliação de política):} qual é o valor de uma política dada?
  \item \textbf{Controle:} qual é a melhor política possível?
\end{enumerate}

A pergunta que conduz a aula: \emph{é possível construir algoritmos que, a cada iteração
adicional de computação, melhorem a política monotonicamente? Todos os algoritmos têm
essa propriedade?} A resposta, como veremos, é ``sim, é possível --- mas não é o que
acontece com todos''.

%% ============================================================
\section{Processos de recompensa de Markov}
%% ============================================================

\begin{defbox}{propriedade de Markov}
O futuro é independente do passado dado o presente:
$\Prob(s_{t+1} \mid s_t, a_t, s_{t-1}, a_{t-1}, \dots) = \Prob(s_{t+1} \mid s_t, a_t)$.
O estado é uma \emph{estatística suficiente} da história.
\end{defbox}

Markov não é uma propriedade do mundo, e sim da \emph{representação de estado} que
escolhemos: se a posição do rover não basta para prever o próximo passo, inclua a
velocidade no estado --- e a propriedade volta a valer. A modelagem se organiza em uma
hierarquia em que cada nível acrescenta exatamente um ingrediente:

\begin{itemize}
  \item \textbf{Processo de Markov} $(\gS, P)$: apenas aleatoriedade;
  \item \textbf{Processo de recompensa de Markov} (MRP) $(\gS, P, \rec, \gamma)$:
        acrescenta recompensa --- já é possível \emph{avaliar};
  \item \textbf{Processo de decisão de Markov} (MDP) $(\gS, \gA, P, \rec, \gamma)$:
        acrescenta ação --- agora há o que \emph{decidir}.
\end{itemize}

\begin{defbox}{horizonte, retorno e função valor}
O \textbf{horizonte} $H$ é o número de passos de tempo por episódio (possivelmente
infinito). O \textbf{retorno} é a soma descontada das recompensas a partir de $t$,
\[
  G_t = r_t + \gamma\, r_{t+1} + \gamma^2 r_{t+2} + \cdots + \gamma^{H-1} r_{t+H-1},
\]
e a \textbf{função valor de estado} é o retorno esperado partindo de $s$:
$V(s) = \E[\, G_t \mid s_t = s\,]$.
\end{defbox}

\begin{ideiabox}
$G_t$ é uma variável \emph{aleatória} --- muda a cada trajetória sorteada. $V(s)$ é um
\emph{número}: a esperança já eliminou o acaso. Boa parte da aprendizagem por reforço
consiste em estimar $V$ a partir de amostras de $G_t$.
\end{ideiabox}

\subsection{Por que descontar o futuro?}

Há quatro razões, de naturezas distintas. \textbf{Matemática:} com recompensas limitadas,
$\abs{r} \le \Rmax$, e $\gamma < 1$, a série geométrica converge e
$\abs{G_t} \le \Rmax/(1-\gamma)$ --- sem isso, o valor pode divergir e comparar políticas
perde o sentido. \textbf{Modelagem:} $\gamma$ pode ser lido como a probabilidade de o
episódio \emph{continuar} a cada passo. \textbf{Econômica:} recompensa hoje vale mais do
que amanhã. \textbf{Algorítmica:} é $\gamma<1$ que torna o operador de Bellman uma
\emph{contração} --- a fonte de todas as garantias de convergência desta aula.

Uma regra prática útil: o \emph{horizonte efetivo} de um problema descontado é da ordem
de $1/(1-\gamma)$ passos ($\gamma=0{,}9 \Rightarrow{} \sim 10$ passos;
$\gamma = 0{,}99 \Rightarrow{} \sim 100$).

\begin{alertabox}
$\gamma = 1$ só é seguro em horizonte finito, ou quando um estado terminal é atingido
com probabilidade $1$ (tarefa episódica própria). Fora dessas condições, $V$ pode
divergir.
\end{alertabox}

\subsection{Equação de Bellman e suas soluções}

A propriedade de Markov permite quebrar o retorno em \emph{agora} $+$ \emph{depois}:
\[
  V(s) \;=\; \termo{\rec(s)}{recompensa imediata}
  \;+\; \termo{\gamma \sum_{s' \in \gS} \tran{s'}{s}\, V(s')}{soma descontada do futuro}
\]
Todo o valor de uma trajetória infinita cabe em uma equação com um único passo de
antecipação --- é o princípio de otimalidade de Bellman. Atenção: isso é um
\emph{sistema} de $N = \abs{\gS}$ equações acopladas, uma por estado, e não uma fórmula
fechada. Há duas maneiras clássicas de resolvê-lo.

\paragraph{Solução analítica.} Empilhando os estados em vetores, $V, \rec \in \R^N$ e
$P \in \R^{N\times N}$, a equação vira $V = \rec + \gamma P V$, donde
\[
  (I - \gamma P)\, V = \rec
  \qquad\Longrightarrow\qquad
  \mdestaque{rlLaranja}{V = (I - \gamma P)^{-1} \rec}
\]

\begin{teobox}{a inversa sempre existe se $\gamma < 1$}
Como $\norminf{\gamma P} = \gamma < 1$, a série de Neumann converge e
$(I - \gamma P)^{-1} = \sum_{k \ge 0} (\gamma P)^k$. O termo de índice $k$ tem leitura
própria: é a contribuição das recompensas $k$ passos à frente.
\end{teobox}

O custo é uma inversão de matriz, $\Theta(N^3)$ --- viável apenas para $N$ moderado.

\paragraph{Solução iterativa (programação dinâmica).}

\begin{algbox}{avaliação iterativa de um MRP}
Inicialize $V_0(s) = 0$ para todo $s$. Para $k = 1, 2, \dots$ até convergir
(por exemplo, até $\norminf{V_k - V_{k-1}} \le \varepsilon$), atualize, para todo
$s \in \gS$:
\[
  V_k(s) \leftarrow \rec(s) + \gamma \sum_{s' \in \gS} \tran{s'}{s}\, V_{k-1}(s')
\]
Custo por varredura: $O(\abs{\gS}^2)$.
\end{algbox}

\begin{ideiabox}
$V_k(s)$ tem uma leitura própria: é o valor \emph{exato} de $s$ quando restam exatamente
$k$ passos de episódio. E na comparação de custos --- analítico $\Theta(N^3)$ uma única
vez \emph{versus} iterativo $O(N^2)$ por varredura --- o iterativo vence para $N$ grande.
Mais importante: é ele que sobrevive à passagem para o caso sem modelo.
\end{ideiabox}

%% ============================================================
\section{Processos de decisão de Markov e políticas}
%% ============================================================

\begin{defbox}{processo de decisão de Markov}
Um MDP é a tupla $(\gS, \gA, P, \rec, \gamma)$, em que $\gS$ é um conjunto finito de
estados de Markov; $\gA$ é um conjunto finito de ações; $P$ é o modelo de
dinâmica/transição, um por ação, $\tran{s_{t+1}=s'}{s_t=s,\, a_t=a}$; $\rec$ é a função
de recompensa, $\rec(s,a) = \E[\, r_t \mid s_t = s,\, a_t = a\,]$; e $\gamma \in [0,1]$
é o fator de desconto. Em resumo: MDP $=$ MRP $+$ ações.
\end{defbox}

\begin{alertabox}
A recompensa aparece na literatura em três formatos --- $\rec(s)$, $\rec(s,a)$ e
$\rec(s,a,s')$. Neste curso usaremos predominantemente $\rec(s,a)$. Os três têm o mesmo
poder expressivo (basta tomar esperanças), mas trocar de convenção no meio de uma
dedução é uma fonte clássica de erro de sinal e de índice.
\end{alertabox}

\paragraph{Exemplo condutor: o \emph{Mars Rover}.} Sete estados ($s_1, \dots, s_7$,
a posição do rover em uma linha), duas ações determinísticas ($a_1$: esquerda; $a_2$:
direita; nas bordas, a ação que sairia do mapa mantém o rover parado) e recompensas
$+1$ em $s_1$, $+10$ em $s_7$ e $0$ nos demais estados. Cada linha das matrizes
$P(\cdot \mid \cdot, a_1)$ e $P(\cdot \mid \cdot, a_2)$ é uma distribuição sobre o
próximo estado; por serem determinísticas, cada linha tem um único $1$.

\begin{center}
\begin{tikzpicture}[node distance=5.4mm]
  \foreach \i [count=\c] in {1,...,7}{%
    \ifnum\c=1 \node[estado] (s\i) {$s_{\i}$};
    \else \pgfmathtruncatemacro{\prev}{\i-1}
      \node[estado, right=of s\prev] (s\i) {$s_{\i}$};
    \fi}
  \node[estado destacado] at (s1) {$s_1$};
  \node[estado destacado] at (s7) {$s_7$};
  \foreach \i in {1,...,6}{%
    \pgfmathtruncatemacro{\nxt}{\i+1}
    \draw[trans destac] (s\i) to[bend left=40] (s\nxt);}
  \foreach \i in {2,...,7}{%
    \pgfmathtruncatemacro{\prv}{\i-1}
    \draw[trans] (s\i) to[bend left=40] (s\prv);}
  \node[font=\scriptsize\bfseries, color=rlLaranja, above=3.2mm of s4] {$a_2$: direita};
  \node[font=\scriptsize\bfseries, color=rlCinza,   below=3.2mm of s4] {$a_1$: esquerda};
  \node[recomp, above=3.2mm of s1] {$+1$};
  \node[recomp, above=3.2mm of s7] {$+10$};
\end{tikzpicture}
\end{center}

\begin{defbox}{política}
Uma política especifica que ação tomar em cada estado. Na forma mais geral, é uma
distribuição condicional: $\pol(a \mid s) = \Prob(a_t = a \mid s_t = s)$. Casos
particulares: \textbf{determinística} ($\pol(s) = a$, massa concentrada em uma ação) e
\textbf{estacionária} (independente de $t$).
\end{defbox}

Trabalhar com a forma estocástica desde o início não é preciosismo: ela permite
exploração, é diferenciável nos parâmetros (o que habilitará os métodos de gradiente de
política) e é a única forma ótima em jogos e em MDPs parcialmente observáveis. E uma
distinção que evita confusões futuras: uma política \emph{não} é um plano. Um plano é
uma sequência fixa de ações; uma política é uma \emph{função reativa} do estado --- ela
responde ao que de fato aconteceu, e por isso é robusta à estocasticidade do ambiente.

\subsection{A observação estrutural central: MDP + política = MRP}

Fixar uma política elimina a decisão: sobram apenas aleatoriedade e recompensa. Um MDP
com $\pol$ fixada é exatamente o MRP $(\gS, \rec^{\pol}, P^{\pol}, \gamma)$, com
\[
  \rec^{\pol}(s) = \sum_{a \in \gA} \pol(a \mid s)\, \rec(s,a),
  \qquad
  P^{\pol}(s' \mid s) = \sum_{a \in \gA} \pol(a \mid s)\, \tran{s'}{s,a}.
\]

\begin{ideiabox}
Esta é a observação mais rentável da aula: \emph{tudo} que sabemos calcular para MRPs
--- forma matricial, solução analítica $(I - \gamma P^{\pol})^{-1} \rec^{\pol}$,
algoritmo iterativo --- aplica-se imediatamente à \emph{avaliação} de uma política em um
MDP. O problema de predição não exige teoria nova. O que não vem de graça é o
\emph{controle}: escolher $\pol$. É aí que o operador $\max$ entra em cena, e com ele a
não linearidade.
\end{ideiabox}

%% ============================================================
\section{Avaliação de política (predição)}
%% ============================================================

\begin{algbox}{avaliação iterativa de política}
Inicialize $V_0^{\pol}(s) = 0$. Para $k = 1, 2, \dots$ até convergir, atualize, para
todo $s \in \gS$:
\[
  V_k^{\pol}(s) \leftarrow \sum_{a \in \gA} \pol(a \mid s)
  \Bigl[\, \rec(s,a) + \gamma \sum_{s' \in \gS} \tran{s'}{s,a}\, V_{k-1}^{\pol}(s') \,\Bigr]
\]
Este é o \emph{backup de Bellman para uma política específica}. Se $\pol$ é
determinística, a soma externa desaparece:
$V_k^{\pol}(s) \leftarrow \rec(s, \pol(s)) + \gamma \sum_{s'} \tran{s'}{s,\pol(s)}\, V_{k-1}^{\pol}(s')$.
\end{algbox}

Note o que \emph{não} há aqui: um $\max$. A política já foi fixada, o operador é
\emph{linear} em $V$ --- e é isso que permite, alternativamente, resolver o sistema de
forma fechada.

\begin{exbox}{L2E1 --- uma iteração de avaliação no \emph{Mars Rover}}
Suponha agora que $a_1$ seja \emph{estocástica}: $p(s_6 \mid s_6, a_1) = 0{,}5$ e
$p(s_7 \mid s_6, a_1) = 0{,}5$ (e analogamente nos demais estados). Recompensa: para
toda ação, $+1$ em $s_1$, $+10$ em $s_7$, $0$ nos demais. Seja $\pol(s) = a_1$ para todo
$s$, com $V_k = [\,1\;\,0\;\,0\;\,0\;\,0\;\,0\;\,10\,]$, $k = 1$ e $\gamma = 0{,}5$.
Calcule $V_{k+1}(s_6)$.
\end{exbox}

\paragraph{Resolução.} Aplicamos o backup para política determinística em um único
estado. Como $p(\cdot \mid s_6, a_1)$ tem massa apenas em $s_6$ e $s_7$, somente dois
termos são não nulos:
\[
  V_{k+1}(s_6)
  = \underbrace{0}_{\rec(s_6, a_1)}
  + \underbrace{0{,}5}_{\gamma} \cdot \bigl(\, 0{,}5 \cdot \underbrace{0}_{V_k(s_6)}
  + 0{,}5 \cdot \underbrace{10}_{V_k(s_7)} \,\bigr)
  = \mdestaque{rlVerde}{2{,}5}
\]

\begin{ideiabox}
Repare no mecanismo: o valor \emph{difunde} para trás pelas transições. Na próxima
varredura, $s_5$ também deixará de valer zero. Cada rodada propaga informação um passo
--- por isso o número de iterações necessárias depende do ``diâmetro'' do MDP.
\end{ideiabox}

\begin{discbox}
No \emph{Mars Rover} --- 7 estados, 2 ações --- quantas políticas \emph{determinísticas}
existem? E a política ótima de um MDP (a de maior valor) é única?
\smallskip

\emph{Direção da resposta:} $2^7 = 128$; em geral, $\abs{\gA}^{\abs{\gS}}$, uma escolha
independente por estado. E não: pode haver duas (ou infinitas) políticas com a mesma
função valor máxima --- basta um empate no máximo de $Q$ em algum estado. O que é único
é a \emph{função valor ótima} $\Vstar$, não o argumento que a atinge. Essa distinção
reaparece o curso inteiro.
\end{discbox}

%% ============================================================
\section{Controle por iteração de política}
%% ============================================================

O problema de controle pede $\polstar(s) = \argmax_{\pol} \Vpi(s)$. Para um MDP de
horizonte infinito com $\gamma < 1$, a política ótima é \textbf{determinística}
(aleatorizar não ajuda), \textbf{estacionária} (não depende de $t$) e
\textbf{não necessariamente única} (pode haver empates); já a função valor ótima
$\Vstar$ é única. Leia o $\argmax$ com cuidado: ele afirma que existe uma política que
maximiza o valor em \emph{todos os estados ao mesmo tempo}. Em otimização multiobjetivo
isso em geral não acontece; aqui acontece, por causa da estrutura de Bellman.

Enumerar candidatas está fora de cogitação: há $\abs{\gA}^{\abs{\gS}}$ políticas
determinísticas ($\abs{\gS}=100$ e $\abs{\gA}=4$ já dão $4^{100} \approx 10^{60}$). A
iteração de política faz muito melhor: usa o valor da política atual para construir
diretamente uma política melhor.

\begin{defbox}{valor de estado-ação $Q$}
\[
  Q^{\pol}(s,a) = \rec(s,a) + \gamma \sum_{s' \in \gS} \tran{s'}{s,a}\, V^{\pol}(s')
\]
Leitura: ``tome a ação $a$ \emph{agora} e siga $\pol$ \emph{daí em diante}''. Vale a
identidade $\Vpi(s) = \sum_a \pol(a \mid s)\, Q^{\pol}(s,a)$.
\end{defbox}

$V$ responde ``quão bom é este estado?''; $Q$ responde ``quão boa é esta \emph{ação}
neste estado?'' --- que é a pergunta que se precisa responder para \emph{agir}. Daí o
domínio de $Q$ nos algoritmos sem modelo.

\begin{algbox}{Iteração de Política (PI)}
$i \leftarrow 0$; inicialize $\pol_0(s)$ aleatoriamente. Enquanto $i = 0$ ou
$\norm{\pol_i - \pol_{i-1}}_1 > 0$ (a norma $L_1$ só verifica se a política mudou em
\emph{algum} estado):
\begin{enumerate}
  \item \textbf{Avaliação:} $V^{\pol_i} \leftarrow$ valor de $\pol_i$
        (pela redução MDP${}+\pol{}=\,$MRP da Seção 2);
  \item \textbf{Melhoria:} calcule $Q^{\pol_i}(s,a)$ para todos $s, a$ e tome a política
        gulosa \(\;\pol_{i+1}(s) = \argmax_a Q^{\pol_i}(s,a)\); faça $i \leftarrow i+1$.
\end{enumerate}
\end{algbox}

\begin{ideiabox}
Nada garante \emph{a priori} que $\pol_{i+1}$ seja melhor: ela foi construída com
$V^{\pol_i}$, o valor da política \emph{antiga}. Provar que a troca compensa é o teorema
a seguir. E o esquema geral --- alternar avaliação e melhoria --- é a \emph{iteração de
política generalizada}, esqueleto de quase todo algoritmo do curso, inclusive Q-learning
e métodos ator-crítico.
\end{ideiabox}

\subsection{Melhoria monotônica}

Escrevemos $V^{\pol_1} \ge V^{\pol_2}$ quando $V^{\pol_1}(s) \ge V^{\pol_2}(s)$ para
\emph{todo} $s \in \gS$ (uma ordem parcial entre políticas).

\begin{teobox}{Proposição --- melhoria monotônica}
Seja $\pol_{i+1}$ a política obtida pelo passo de melhoria aplicado a $\pol_i$. Então
$V^{\pol_{i+1}} \ge V^{\pol_i}$, com desigualdade \emph{estrita} em ao menos um estado
se $\pol_i$ for subótima.
\end{teobox}

A intuição em duas frases. Para qualquer $s$, escolher a melhor ação é ao menos tão bom
quanto seguir a política antiga:
$\max_a Q^{\pol_i}(s,a) \ge Q^{\pol_i}(s, \pol_i(s)) = V^{\pol_i}(s)$. Isso mostra que
tomar $\pol_{i+1}$ \emph{por um passo} e depois voltar a $\pol_i$ não piora nada --- mas
a política proposta é seguir $\pol_{i+1}$ \emph{sempre}. O argumento de um passo não
fecha a prova: é preciso aplicá-lo repetidamente, em um telescópio de desigualdades.

\begin{provabox}
\begin{align*}
  V^{\pol_i}(s)
    &\le \max_{a} Q^{\pol_i}(s,a)
     = \max_{a} \Bigl[\, \rec(s,a) + \gamma \sum_{s'} \tran{s'}{s,a}\, V^{\pol_i}(s') \,\Bigr]
     && \text{\color{rlCinza}(i) o máximo domina} \\
    &= \rec(s, \pol_{i+1}(s)) + \gamma \sum_{s'} \tran{s'}{s,\pol_{i+1}(s)}\, V^{\pol_i}(s')
     && \text{\color{rlCinza}(ii) definição de $\pol_{i+1}$} \\
    &\le \rec(s, \pol_{i+1}(s)) + \gamma \sum_{s'} \tran{s'}{s,\pol_{i+1}(s)}
      \Bigl[\, \max_{a'} Q^{\pol_i}(s',a') \,\Bigr]
     && \text{\color{rlCinza}(iii) aplica (i) em $s'$}
\end{align*}
Expandindo o colchete com (ii) aplicado em $s'$, aparece mais um passo executado sob
$\pol_{i+1}$, com a cauda ainda avaliada sob $\pol_i$:
\[
  \le \rec(s, \pol_{i+1}(s)) + \gamma \sum_{s'} \tran{s'}{s,\pol_{i+1}(s)}
      \Bigl[\, \rec(s', \pol_{i+1}(s')) + \gamma \sum_{s''} \tran{s''}{s',\pol_{i+1}(s')}\,
      V^{\pol_i}(s'') \,\Bigr] \le \cdots
\]
Repetindo indefinidamente, cada linha empurra a troca de política um passo adiante; o
resíduo avaliado sob $\pol_i$ é multiplicado por $\gamma^k \to 0$ e desaparece no
limite, restando exatamente a definição de $V^{\pol_{i+1}}(s)$. Logo
$V^{\pol_i}(s) \le V^{\pol_{i+1}}(s)$ para todo $s$. \hfill $\blacksquare$
\end{provabox}

\begin{discbox}
Duas perguntas sobre o comportamento de PI. (1) Se a política não muda em uma rodada,
ela pode voltar a mudar depois? (2) Existe um número máximo de rodadas?
\smallskip

\emph{Direção da resposta:} (1) Não. Se $\pol_{i+1} = \pol_i$, então
$Q^{\pol_{i+1}} = Q^{\pol_i}$ (mesma política, mesmo valor), logo
$\pol_{i+2} = \argmax_a Q^{\pol_{i+1}} = \argmax_a Q^{\pol_i} = \pol_{i+1}$: ponto fixo
--- e esse ponto fixo é ótimo. (2) Sim: $\abs{\gA}^{\abs{\gS}}$. Pela monotonia, cada
política só aparece em uma rodada (a menos que já seja ótima). Na prática, PI converge
em poucas dezenas de rodadas; Ye (2011) mostrou que, para $\gamma$ fixo, PI é
\emph{fortemente polinomial} em $\abs{\gS}$ e $\abs{\gA}$.
\end{discbox}

%% ============================================================
\section{Controle por iteração de valor}
%% ============================================================

A segunda estratégia de controle não mantém uma política: mantém o \emph{valor ótimo}
de começar em cada estado supondo que restam $k$ passos de episódio --- e itera
considerando episódios cada vez mais longos. São dois eixos de aproximação ortogonais:
PI é exata no \emph{valor} e aproxima a \emph{política}; VI é exata na otimalidade local
e aproxima o \emph{horizonte}. Ambas convergem para o mesmo par $(\Vstar, \polstar)$.

\begin{defbox}{operador de backup de Bellman $\Bop$}
Aplicado a uma função valor qualquer, devolve uma nova função valor sobre todos os
estados, melhorando o valor quando há melhora possível:
\[
  \Bop V(s) = \max_{a} \Bigl[\, \rec(s,a) + \gamma \sum_{s' \in \gS} \tran{s'}{s,a}\, V(s') \,\Bigr]
\]
\end{defbox}

A troca da \emph{soma ponderada por $\pol$} pelo \emph{$\max$ sobre ações} é o que
distingue predição de controle --- e o que destrói a linearidade: com $\Bop$, a solução
fechada por inversão de matriz deixa de existir.

\begin{algbox}{Iteração de Valor (VI)}
Inicialize $V_0(s) = 0$. Repita $V_{k+1} = \Bop V_k$ até convergir (por exemplo,
$\norminf{V_{k+1} - V_k} \le \epsilon$). Ao final, extraia a política \emph{gulosa} em
relação ao último valor:
$\pol(s) = \argmax_a \bigl[\, \rec(s,a) + \gamma \sum_{s'} \tran{s'}{s,a}\, V_{k+1}(s') \,\bigr]$.
\end{algbox}

Na linguagem de operadores, as duas famílias se unificam. Definindo o backup para uma
política fixa, $\Bpol V(s) = \rec^{\pol}(s) + \gamma \sum_{s'} P^{\pol}(s' \mid s) V(s')$,
a avaliação de política é o cálculo do \emph{ponto fixo} de $\Bpol$
($\Vpi = \Bpol \Bpol \cdots \Bpol V$), e a melhoria é um único backup com $\max$ seguido
de $\argmax$. VI itera $\Bop$ (com $\max$) desde o início; PI itera $\Bpol$ (linear) até
o ponto fixo e só então usa o $\max$ uma vez.

\begin{alertabox}
Em VI, os valores intermediários $V_k$ em geral \emph{não} correspondem ao valor de
nenhuma política executável --- são valores ótimos de um problema truncado em $k$
passos. Consequência: a qualidade das políticas gulosas extraídas ao longo do caminho
pode oscilar antes de estabilizar. A monotonia rodada a rodada é uma garantia de PI, não
de VI.
\end{alertabox}

\subsection{Por que VI converge: contração}

\begin{defbox}{operador de contração}
Seja $O$ um operador e $\norm{\cdot}$ uma norma. Se $\norm{OV - OV'} \le \norm{V - V'}$
para todos $V, V'$, então $O$ é um operador de \emph{contração} (estrita, se a
desigualdade vale com um fator $\gamma < 1$).
\end{defbox}

Imagem mental: um mapa que \emph{aproxima os pontos}. Aplicado repetidamente, todos os
pontos colapsam em um único ponto fixo --- não importa de onde se parta.

\begin{teobox}{o backup de Bellman é uma contração para $\gamma<1$}
Na norma do supremo, $\norminf{V - V'} = \max_s \abs{V(s) - V'(s)}$, vale
$\norminf{\Bop V - \Bop V'} \le \gamma\, \norminf{V - V'}$.
\end{teobox}

\begin{provabox}
Para cada estado $s$, usando primeiro que a diferença de máximos é limitada pelo máximo
das diferenças, $\abs{\max_a f(a) - \max_{a'} g(a')} \le \max_a \abs{f(a) - g(a)}$:
\begin{align*}
  \abs{\Bop V_k(s) - \Bop V_j(s)}
    &\le \max_a \Bigl|\, \rec(s,a) + \gamma \sum_{s'} \tran{s'}{s,a}\, V_k(s')
       - \rec(s,a) - \gamma \sum_{s'} \tran{s'}{s,a}\, V_j(s') \,\Bigr| \\
    &= \max_a \, \gamma \Bigl|\, \sum_{s'} \tran{s'}{s,a} \bigl( V_k(s') - V_j(s') \bigr) \,\Bigr|
     \;\le\; \max_a \, \gamma \sum_{s'} \tran{s'}{s,a}\, \norminf{V_k - V_j} \\
    &= \gamma\, \norminf{V_k - V_j} \qquad \text{\color{rlCinza}(as probabilidades somam 1).}
\end{align*}
Tomando o máximo sobre $s$ obtém-se a contração. Note dois detalhes: a recompensa
$\rec$ \emph{cancela} --- a contração não depende dela ---, e mesmo que todas as
desigualdades fossem igualdades ainda haveria contração, desde que $\gamma<1$.
\hfill $\blacksquare$
\end{provabox}

\begin{teobox}{consequências (ponto fixo de Banach)}
Se $\gamma < 1$: (i) $\Bop$ tem um \emph{único} ponto fixo, que é $\Vstar$; (ii) VI
converge a partir de \emph{qualquer} inicialização $V_0$; (iii) a convergência é
geométrica, $\norminf{V_k - \Vstar} \le \gamma^k \norminf{V_0 - \Vstar}$. Além disso, se
o critério de parada $\norminf{V_{k+1} - V_k} \le \epsilon$ foi atingido, então
$\norminf{V_{k+1} - \Vstar} \le \gamma\epsilon/(1-\gamma)$, e a política gulosa extraída
perde no máximo $2\gamma\epsilon/(1-\gamma)$ em valor.
\end{teobox}

Duas consequências práticas: a inicialização não afeta o destino, só o tempo de viagem
(e uma boa inicialização acelera); e $\gamma$ perto de $1$ torna a convergência lenta
--- o número de iterações cresce como $1/(1-\gamma)$.

%% ============================================================
\section{Horizonte finito e avaliação por simulação}
%% ============================================================

\subsection{VI em horizonte finito}

Em horizonte finito, VI deixa de ser aproximação e vira \emph{o algoritmo exato}:
$V_k$ e $\pol_k$ são o valor e a política ótimos \emph{quando restam $k$ decisões}.
Basta rodar o laço de VI por exatamente $H$ iterações, guardando a cada passo também a
política gulosa $\pol_{k+1}(s) = \argmax_a [\, \rec(s,a) + \gamma \sum_{s'} \tran{s'}{s,a} V_k(s') \,]$.
Para agir no instante $t$, usa-se $\pol_{H-t}$ --- a política indexada por \emph{quantos
passos restam}.

\begin{discbox}
Em tarefas de horizonte \emph{finito}, a política ótima é estacionária (independente do
passo de tempo)?
\smallskip

\emph{Direção da resposta:} em geral, não --- a ação ótima depende de quanto tempo
resta. No \emph{Mars Rover}, em $s_4$ com muitos passos restantes vale caminhar até o
$+10$ em $s_7$; com um único passo restante, nenhuma direção alcança recompensa; com
dois ou três, o $+1$ de $s_1$ e o $+10$ de $s_7$ competem de forma diferente a cada
instante. A política ótima precisa ser indexada pelo tempo: $\pol_t^*(s)$.
\end{discbox}

\subsection{Avaliação por simulação (Monte Carlo)}

Alternativa para estimar o valor de uma política: gere um grande número de
\emph{episódios} seguindo-a, calcule o retorno de cada um e tire a \emph{média}.
Desigualdades de concentração (como a de Hoeffding) limitam a rapidez com que a média se
aproxima da esperança. Nenhuma suposição markoviana é necessária --- e, mais importante,
nenhum \emph{modelo}: basta poder interagir. O preço é a variância: cada retorno é uma
amostra ruidosa, e a precisão melhora apenas com $1/\sqrt{n}$ episódios.

\paragraph{Exemplo numérico.} No \emph{Mars Rover} estocástico (passeio aleatório: de
cada estado interior, prob.\ $0{,}4$ para cada vizinho e $0{,}2$ de ficar parado;
$0{,}6$ de ficar parado nas bordas), com episódios de $H = 4$ passos a partir de $s_4$ e
$\gamma = \tfrac12$:

\begin{center}
\begin{tabular}{@{}llc@{}}
\toprule
\textbf{Episódio} & \textbf{Retorno } $G = r_1 + \tfrac12 r_2 + \tfrac14 r_3 + \tfrac18 r_4$ & \textbf{Valor} \\
\midrule
$s_4, s_5, s_6, s_7$ & $0 + \tfrac12 \cdot 0 + \tfrac14 \cdot 0 + \tfrac18 \cdot 10$ & $1{,}25$ \\
$s_4, s_4, s_5, s_4$ & $0 + \tfrac12 \cdot 0 + \tfrac14 \cdot 0 + \tfrac18 \cdot 0$  & $0$ \\
$s_4, s_3, s_2, s_1$ & $0 + \tfrac12 \cdot 0 + \tfrac14 \cdot 0 + \tfrac18 \cdot 1$  & $0{,}125$ \\
\bottomrule
\end{tabular}
\end{center}

Três amostras, três respostas --- todas legítimas. $V(s_4)$ é a média dessa distribuição
(o cálculo exato por programação dinâmica dá $V(s_4) \approx 0{,}088$); estimá-la bem
exige muitos episódios. DP calcula a média sem nunca sortear; Monte Carlo sorteia sem
nunca conhecer $P$. A Aula 3 vive exatamente nesse contraste.

%% ============================================================
\section{Síntese}
%% ============================================================

\subsection{Comparação dos métodos}

\paragraph{Três formas de avaliar uma política.}

\begin{center}
\begin{tabular}{@{}p{3.2cm}p{3.6cm}cc@{}}
\toprule
\textbf{Método} & \textbf{Custo} & \textbf{Exige Markov?} & \textbf{Exige modelo?} \\
\midrule
Analítico $(I - \gamma P^{\pol})^{-1} \rec^{\pol}$ & $\Theta(N^3)$, uma vez & sim & sim \\
Iterativo (DP) & $O(N^2)$ por varredura & sim & sim \\
Simulação (Monte Carlo) & $O(\text{episódios} \times H)$ & \textbf{não} & não (só amostras) \\
\bottomrule
\end{tabular}
\end{center}

\paragraph{Os dois algoritmos de controle.}

\begin{center}
\begin{tabular}{@{}p{4.0cm}p{4.5cm}p{4.5cm}@{}}
\toprule
 & \textbf{Iteração de valor} & \textbf{Iteração de política} \\
\midrule
O que mantém & valor ótimo p/ horizonte $k$ & política + seu valor $\infty$ \\
Passo & $V_{k+1} = \Bop V_k$ & avaliar $\Bpol$; melhorar c/ $\max$ \\
Custo por iteração & $O(\abs{\gS}^2 \abs{\gA})$ & avaliação $+$ $O(\abs{\gS}^2 \abs{\gA})$ \\
N\textsuperscript{o} de iterações & muitas e baratas & poucas e caras \\
Monotonia da política & não garantida & \textbf{garantida} \\
Parente em RL moderno & Q\mbox{-}learning & gradiente de política \\
\bottomrule
\end{tabular}
\end{center}

\begin{ideiabox}
\emph{Iteração de política modificada}: no passo de avaliação, em vez de iterar $\Bpol$
até o ponto fixo, aplique-o apenas $m$ vezes. Com $m = 1$ recupera-se VI; com
$m = \infty$, PI. Quase tudo entre os dois extremos funciona --- e costuma ser mais
rápido que ambos.
\end{ideiabox}

\subsection{Vocabulário}

\textbf{Modelo}: representação matemática da dinâmica ($P$) e da recompensa ($\rec$).
\textbf{Política}: função que mapeia estados em ações (ou distribuições sobre ações).
\textbf{Função valor}: recompensas futuras esperadas a partir de um estado (e/ou ação),
seguindo uma política específica. Estes três objetos organizam todo o campo: algoritmos
baseados em modelo, em valor ou em política --- e suas combinações.

\subsection{O que você deve saber ao final desta aula}

\begin{itemize}
  \item \textbf{Definir}: MP, MRP, MDP, operador de Bellman, contração, modelo, valor
        $Q$, política.
  \item \textbf{Implementar}: iteração de valor e iteração de política.
  \item \textbf{Comparar}: prós e contras das abordagens de avaliação de política.
  \item \textbf{Provar}: propriedades de contração e a melhoria monotônica.
  \item \textbf{Criticar}: limitações das abordagens apresentadas e da hipótese de
        Markov --- em particular, saber dizer quais métodos de avaliação exigem a
        hipótese de Markov e quais não.
\end{itemize}

\subsection{Prática recomendada}

\begin{exbox}{--- para consolidar fora de aula}
\begin{enumerate}
  \item Prove que a iteração de valor converge para uma solução \emph{única} em espaços
        discretos de estados e ações com $\gamma < 1$.
  \item A inicialização dos valores em VI afeta alguma coisa? O quê, exatamente --- e o
        que ela \emph{não} afeta?
  \item O valor da política extraída de VI a cada rodada é garantidamente monotônico
        (se executada no problema real de horizonte infinito), como na iteração de
        política? \emph{(Esta é a pergunta mais sutil da aula --- compare com a caixa
        de atenção da Seção 5.)}
  \item Implemente VI e PI para o \emph{Mars Rover} e meça: quantas iterações cada um
        leva com $\gamma = 0{,}5$? E com $\gamma = 0{,}95$? O laboratório interativo da
        aula permite conferir suas respostas.
\end{enumerate}
\end{exbox}

\vspace{1em}
\begin{center}
{\color{rlCinza}\small\itshape Próxima aula: avaliação de política sem modelo ---
Monte Carlo e diferenças temporais.}
\end{center}

\end{document}
